\documentclass[../../main.tex]{subfiles}
\graphicspath{{\subfix{../../image/}}}

\begin{document}

\section*{第五章小结}\label{section:ch05-小结}
\addcontentsline{toc}{section}{第五章小结}

本章研究了极限定理，这是概率论基础中比较深入的结果；前几章学到的知识在这里得到了综合应用；一些重要问题在这里进一步讨论并获得解决。

在我们的课程中，为了使读者对极限定理有直观的认识，是从 Bernoulli 试验场合开始叙述的。这里所用的工具比较初等：Bernoulli 大数定律是用矩法证明的；De Moivre--Laplace 定理则通过利用 Stirling 公式进行渐近估计而得到。

接着我们处理独立同分布场合，这是 Bernoulli 试验的直接推广，也是在实际中，特别是数理统计中，最常碰到的情况。为了证明 Khinchin 大数定律及 Lindeberg--Lévy 定理已用到特征函数，所用的方法具有普遍性及简明性，是读者比较容易理解的。应当指出，收敛性概念及特征函数的连续性定理是深入研究极限理论所不可缺少的，所以对这部分内容我们预先作了相当详细的叙述。后面用到的主要是结论，因此有关证明在初学时不妨略去。

最后，介绍了强大数定律及一般场合的中心极限定理，这是概率论中相当深刻的结果。前者的证明通过建立比 Chebyshev 不等式更为锐利的不等式而实现；后者的证明则得力于特征函数这一有力的工具的巧妙应用。到此为止，概率论中提出的古典极限定理问题已获得了令人满意的解决。

本章对极限定理的处理采用模块式结构，即分 Bernoulli 试验、独立同分布、一般等三个场合进行，与历史发展大致平行，这样设计的目的一是体现循序渐进、由浅入深，二是让本书能更好地适应不同教学要求。

纵览整个发展过程，令人印象深刻的是，经典问题的最终解决，主要靠工具的改进，从直接展开，到矩法，再到特征函数法；也靠方法的精密化，例如从 Chebyshev 不等式，到 Markov 不等式，再到 Kolmogorov 等更精细的不等式。

独立同分布模型事实上成为概率论继古典概型、Bernoulli 概型之后的第三个重要概型，直至近代才让位于随机过程。

\end{document}
